\subsection{Finite expected work suffices for the first score identity}
\label{sec:new-finite-score}

The usual first-order score argument needs no exponential tail for a
Bernoulli factory with bounded output.  This observation removes the
extra regularity hypothesis from the independent-recursion obstruction.
It is a statement about sequential experiments, and we make no novelty
claim for the information principle.  We give a proof because the
finite-mean hypothesis is the one used by the sampling problem.

\begin{lemma}[A first derivative under a finite-mean stopping rule]
\label{lem:new-finite-score}
Let $I\subset(-1,1)$ be an open interval containing zero.  For $z\in I$, let
$X_1,X_2,\ldots$ be independent signs of mean $z$.  A fixed randomized
algorithm uses an auxiliary random tape whose law does not depend on
$z$, requests these signs sequentially, and terminates almost surely
for every $z\in I$.  Let $T$ be its number of source requests and let
$Y\in[-1,1]$ be its output.  If $\mathbb E_0T<\infty$, then
$f(z)=\mathbb E_zY$ is differentiable at zero, and
\[
 f'(0)=\mathbb E_0[YS_T],\qquad
 S_T=\sum_{j=1}^T X_j,\qquad
 \mathbb E_0 S_T^2=\mathbb E_0T.
\]
No bound on $\mathbb E_zT$ away from zero is assumed.
\end{lemma}
\paragraph{Proof idea.}
Truncate the stopped transcript at a finite time, where likelihood
differentiation is elementary. A relative-entropy tail bound and the stopped
score isometry let the truncation tend to infinity under only a finite mean
stopping time at the bias of interest.

\begin{proof}
All random tapes may be taken in a standard Borel space, for example
a countable sequence of fair bits.  Include the complete auxiliary
tape in the initial sigma-field.  Conditional on this tape, $T$ is a
stopping time for the source signs.  For an integer $N\ge0$, let
$\mathcal F_N$ be the stopped prefix sigma-field at $T\wedge N$, and
put $g_N=\mathbb E_0[Y\mid\mathcal F_N]$.  The finite-prefix likelihood
ratio under mean $z$, relative to mean zero, is
\[
 L_{z,N}=\prod_{j=1}^{T\wedge N}(1+zX_j)
       \le(1+|z|)^N.
\]
Consequently $f_N(z)=\mathbb E_z g_N$ is a polynomial of degree at
most $N$, with
\[
 f_N(0)=f(0),\qquad
 f_N'(0)=\mathbb E_0[g_NS_{T\wedge N}]
        =\mathbb E_0[YS_{T\wedge N}].
\]
The integration over the auxiliary tape preserves this finite-degree
polynomial identity.

Write $d(z)=-\tfrac12\log(1-z^2)$, the relative entropy of a fair
sign with respect to a sign of mean $z$.  Conditional on a reached
prefix, the relative entropy of the remaining stopped transcript
under mean zero with respect to mean $z$ is
\[
 d(z)\,\mathbb E_0[(T-N)_+\mid\mathcal F_N].
\]
To justify this formula directly, the log likelihood ratio is a sum
of source increments.  Their absolute values are bounded for each
fixed $z\in(-1,1)$, and the remaining expected number of increments
is finite almost surely under mean zero.  Conditional Wald summation
therefore applies.  The auxiliary tape contributes zero relative
entropy.  Almost-sure termination under both parameters makes these
stopped-transcript laws probability measures.

Pinsker's inequality and $|Y|\le1$ bound the conditional change of
its mean by the square root of twice that relative entropy.  Averaging
under the mean-$z$ prefix law and using the displayed bound on
$L_{z,N}$ gives
\begin{align*}
 |f(z)-f_N(z)|
 &\le (1+|z|)^N
      \sqrt{2d(z)}\,
      \mathbb E_0\!\left[
       \sqrt{\mathbb E_0[(T-N)_+\mid\mathcal F_N]}\right]\\
 &\le (1+|z|)^N\sqrt{2d(z)}
           \sqrt{\mathbb E_0(T-N)_+}.
\end{align*}
The second line is Jensen's inequality.  Since
$\sqrt{2d(z)}/|z|\longrightarrow1$, it follows that
\[
 \limsup_{z\to0}
 \left|\frac{f(z)-f(0)}z-f_N'(0)\right|
 \le\sqrt{\mathbb E_0(T-N)_+}.
\]

The fair-sign martingale isometry gives
$\mathbb E_0 S_{T\wedge N}^2=\mathbb E_0(T\wedge N)$.
For $M>N$ it also gives
\[
 \mathbb E_0|S_{T\wedge M}-S_{T\wedge N}|^2
   =\mathbb E_0[(T\wedge M)-(T\wedge N)]
   \le\mathbb E_0(T-N)_+.
\]
Thus the stopped sums converge to $S_T$ in $L^2$, proving the stated
isometry and $f_N'(0)\to\mathbb E_0[YS_T]$.  Letting $N\to\infty$
in the previous limsup bound proves differentiability and the score
identity.
\end{proof}

\begin{corollary}[The integer obstruction requires only finite mean]
\label{cor:new-integer-score}
Under Lemma~\ref{lem:new-finite-score}, retain the assumption $Y\in[-1,1]$.
For every integer $k\ge0$,
\[
 \mathbb E_0T\ge(2k+1)|f'(0)|-k(k+1).
\]
In particular, every finite-mean factory for
$G_r(z)=\tanh(rz)/\tanh r$ obeys
\[
 \mathbb E_0T\ge3r\coth r-2
              =1+r^2+O(r^4).
\]
\end{corollary}
\paragraph{Proof idea.}
The stopped score is integer-valued. The gap between two consecutive integer
magnitudes gives a stronger pointwise inequality than Cauchy--Schwarz;
expectation and the score identity turn it into a source-cost bound.

\begin{proof}
For every integer $v$, the consecutive integers $k$ and $k+1$ give
$(|v|-k)(|v|-k-1)\ge0$.  Hence
\[
 (2k+1)|v|\le v^2+k(k+1).
\]
Apply this to $v=S_T$, use the score identity, and use the stopped
isometry.  The last assertion takes $k=1$ and
$G_r'(0)=r\coth r$.
\end{proof}

For $a=|f'(0)|$ and $k=\lfloor a\rfloor$, the strongest line in
this family is
\[
 (2k+1)a-k(k+1)
 =a^2+(a-k)(k+1-a)\in[a^2,a^2+1/4].
\]
The proof uses only $|Y|\le1$, so postponing the final Bernoulli
rounding of a bounded real estimator does not avoid this local obstruction.
The small amplification $a=1+u$ therefore costs at least $1+3u$,
whereas a single large amplification of size $a$ has the lower bound
$a^2+O(1)$. This distinction explains the loss under repeated fresh
composition; it supplies no upper bound for a composite factory.

\begin{proposition}[A finite-mean independent-recursion obstruction]
\label{prop:new-local-obstruction}
Put $R_0=1$ and $R_{j+1}=\tanh R_j$.  Consider the scalar critical
mean chain, represented by independently invoking a scalar factory
for $G_{R_{j-1}}$ at each layer.  Every source request initiates a
fresh recursive call, and the top normalized factory is entered
with probability $R_D$.  If all local factories terminate almost
surely for all source means and have finite expected work at mean
zero, their expected number of bottom source requests at mean zero
is at least
\[
 e^{-3/4}R_D^{-2}
 \ge e^{-3/4}(1+2D/3)>D/4.
\]
\end{proposition}
\paragraph{Proof idea.}
At mean zero, every local invocation faces the same derivative constraint.
Predictable charging multiplies the resulting local lower bounds through the
chain, preventing a purely local improvement of the critical exponent.

\begin{proof}
Write $a_j=R_{j-1}/R_j$.  Corollary~\ref{cor:new-integer-score}
gives local mean offspring at least $3a_j-2$.  Each conditional
source mean is zero in the mean-zero chain.  Since every source is
fresh, predictable charging and Tonelli's theorem give the bottom
request bound
\[
 R_D\prod_{j=1}^D(3a_j-2).
\]
For $b\ge0$,
\[
 3\log(1+b)-\log(1+3b)
 =\int_0^b\frac{6u}{(1+u)(1+3u)}\,du\le3b^2.
\]
Inequality \eqref{eq:xcothx} gives
$a_j-1\le R_{j-1}^2/3$, and Lemma~\ref{lem:radius} at gain one gives
\[
 R_j^2\le\frac3{2j+3},\qquad R_D^{-2}\ge1+2D/3.
\]
Convexity and midpoint integration give
\[
 \sum_{j\ge0}\frac9{(2j+3)^2}
 \le\int_{-1/2}^{\infty}\frac9{(2t+3)^2}\,dt=\frac94.
\]
It follows that
\begin{align*}
 \prod_{j=1}^D(3a_j-2)
 &\ge\left(\prod_{j=1}^D a_j\right)^3
     \exp\left(-3\sum_{j=1}^D(a_j-1)^2\right)\\
 &\ge R_D^{-3}e^{-3/4}.
\end{align*}
Multiplication by the root gate proves the result.  Finally
$e^{-3/4}(2/3)>1/4$, for example from $e^{3/4}<8/3$.
\end{proof}

The obstruction counts fresh recursive source calls.  It permits every
finite-mean scalar factory, but it does not restrict methods that sample
a composite function directly, share information across layers, or use a
different probabilistic representation.  The global scalar mean-chain
factory already attains order $\sqrt D$.  The general zero-bias
norm-one network exponent remains between $1/2$ and $1$.

\subsection{A dimension-free smoothness certificate}
\label{sec:new-critical-smoothness}

A second-order approach to a composite factory has a useful analytic
bound.  The following proposition concerns derivatives of the real
network extension.  It is not an exact-sampling algorithm.
For a derivative tensor, write $\|\cdot\|_{1,\mathrm{ent}}$ for the
sum of the absolute values of all its entries.

\begin{proposition}[Derivatives of zero-bias contractive tanh networks]
\label{prop:new-critical-derivatives}
Let $h_0(x)=x\in[-1,1]^n$, let
$h_\ell(x)=\tanh(W_\ell h_{\ell-1}(x))$ coordinatewise, and suppose
all rows of every $W_\ell$ have $\ell_1$ norm at most one.
For every output coordinate $f=h_{D,i}$, every $x\in[-1,1]^n$, and
$D\ge1$,
\[
 \|Df\|_{1,\mathrm{ent}}\le1,\quad
 \|D^2f\|_{1,\mathrm{ent}}\le2\sqrt6\sqrt D,\quad
 \|D^3f\|_{1,\mathrm{ent}}\le38D,\quad
 \|D^4f\|_{1,\mathrm{ent}}\le624D^{3/2}.
\]
The derivatives at boundary points are those of the analytic real
extension of the network.
\end{proposition}
\paragraph{Proof idea.}
Differentiate the network through fourth order while retaining the shrinking
activation radius. Even derivatives of tanh carry an extra radius factor,
which improves the usual chain-rule growth estimates.

\begin{proof}
Let $R_0=1$, $R_\ell=\tanh R_{\ell-1}$.  Then
$|h_{\ell,i}|\le R_\ell$ and
$R_\ell\le\sqrt{3/(2\ell+3)}$.
Let $A_\ell,B_\ell,C_\ell,E_\ell$ bound, respectively, the largest
entrywise $\ell_1$ norms of the first through fourth derivative tensors
of a coordinate at depth $\ell$, uniformly over $x$.
At depth zero these bounds are $1,0,0,0$.  Linear row contraction
preserves each bound.  Tensor products multiply entrywise $\ell_1$
norms, and the scalar chain rule therefore applies without a dimension
factor.

Writing $t=\tanh u$, direct differentiation gives
\[
 |\tanh' u|\le1,\quad
 |\tanh''u|\le2|t|,\quad
 |\tanh'''u|\le2,\quad
 |\tanh''''u|\le16|t|.
\]
For the last bound,
$\tanh''''u=t(16-40t^2+24t^4)$, whose bracket has absolute value
at most $16$ on $[-1,1]$.  The chain rule gives $A_\ell\le1$ and
\begin{align*}
 B_\ell&\le B_{\ell-1}+2R_\ell,\\
 C_\ell&\le C_{\ell-1}+6R_\ell B_{\ell-1}+2,\\
 E_\ell&\le E_{\ell-1}+8R_\ell C_{\ell-1}
      +6R_\ell B_{\ell-1}^2+12B_{\ell-1}+16R_\ell.
\end{align*}
The decreasing-function integral bound gives
$\sum_{j=1}^{\ell}R_j\le\sqrt6\sqrt\ell$, and hence
$B_\ell\le2\sqrt6\sqrt\ell$.
Since $R_\ell\sqrt{\ell-1}\le\sqrt{3/2}$, the increment of
$C_\ell$ is at most $36+2=38$.
Using these two bounds, the increment of $E_\ell$ is at most
\[
 \left(448\sqrt{3/2}+24\sqrt6+16\right)\sqrt\ell
 <624\sqrt\ell.
\]
The strict numerical inequality follows from
$\sqrt{3/2}<49/40$ and $\sqrt6<49/20$.
Summing $\sqrt\ell\le\sqrt D$ for $1\le\ell\le D$ completes
the proof.
\end{proof}

A generic factory whose expected work were controlled by the total
second-derivative mass, with access costs charged through the network's
rows, could plausibly use this proposition to reach order $\sqrt D$.
The known tensor Bernstein construction does not supply that interface:
it requests coordinate vectors or evaluates the composite function at
empirical inputs.  Both operations can contain a dimension factor or a
dense network evaluation.  No such factory is proved here.  The
derivative proposition must therefore be kept separate from the
sampling upper bounds.
